% DIN 与 SIM 的对比：问题、历史入口、每个候选的计算量及代价。两列并排，跨栏约148mm。
\begin{tikzpicture}[x=1mm,y=1mm,
  font=\figurefont\footnotesize,
  ctitle/.style={rectangle,rounded corners=1.4pt,minimum height=6.5mm,
    font=\figurefont\bfseries\footnotesize,text=BookInk,align=center,
    text width=74mm,inner sep=2pt},
  cpane/.style={rectangle,draw=BookRule,fill=BookPaper,rounded corners=1.6pt,
    align=left,text=BookInk,text width=66mm,minimum height=54mm,inner sep=4pt}
]
  \node[ctitle,fill=FigureInput!14] (tdin) at (-37,30) {DIN：候选相关激活};
  \node[ctitle,fill=FigureLoss!16] (tsim) at (38,30) {SIM：先检索再精细匹配};

  \node[cpane,anchor=north west] at (-74,27) {
    {\bfseries\textcolor{FigureLoss!85}{解决的问题}}\newline
    一份历史对不同候选应给不同证据；\newline
    先压成单一向量会丢失候选区分。\newline
    {\bfseries\textcolor{FigureLoss!85}{历史入口}}\newline
    候选到来后，对全部 $n$ 条行为逐条激活。\newline
    {\bfseries\textcolor{FigureLoss!85}{每候选计算}}\newline
    $O(n)$ 次候选—行为比较，随历史变长昂贵。\newline
    {\bfseries\textcolor{FigureLoss!85}{收益与代价}}\newline
    表示随候选变化、可解释相关行为；\newline
    长历史下请求时延与算力不可承受。};

  \node[cpane,anchor=north west] at (1,27) {
    {\bfseries\textcolor{FigureLoss!85}{解决的问题}}\newline
    与候选相关的旧兴趣稀疏地埋在很长历史里，\newline
    DIN 无法在预算内逐条扫描。\newline
    {\bfseries\textcolor{FigureLoss!85}{历史入口}}\newline
    GSU 先用类别或内积选 $K\ll n$ 条，\newline
    ESU 只在子集上做候选相关注意力。\newline
    {\bfseries\textcolor{FigureLoss!85}{每候选计算}}\newline
    精细部分降为 $O(K)$；检索依赖预建索引。\newline
    {\bfseries\textcolor{FigureLoss!85}{收益与代价}}\newline
    可覆盖多年历史；但硬截断漏检不可恢复，\newline
    检索与精排相关性需一致（引出 TWIN）。};
\end{tikzpicture}
